\documentclass[10pt,a4paper]{amsart}
\usepackage[margin=19mm]{geometry}
\usepackage{amsmath,amssymb,mathtools}
\usepackage[scr=rsfs,cal=euler]{mathalfa}
\usepackage{xfrac}
\usepackage[unicode,hidelinks,bookmarks=false]{hyperref}
\newcommand{\ie}{i.e.\ }
\newcommand{\cf}{cf.\ }
\newcommand{\resp}{respectively\ }
\newcommand{\Subsection}{\subsection}
\providecommand{\nobreakdash}{\nobreak\hbox{-}}
\setlength{\emergencystretch}{2em}
\multlinegap0pt
\usepackage{amssymb}
\usepackage{enumerate}
\newtheorem{proposition}{Proposition}
\newtheorem{claim}{Claim}
\def\Fo{\mathbf{Fo}}
\def\Formula{\mathbf{A}}
\def\L{\mathbf{L}}
\def\N{\mathbb{N}}
\def\Rule{\mathbf{R}}
\title{Four intuitionistic modal connectives}
\author{Philippe Balbiani, \c{C}igdem Gencer}
\date{Source: arXiv:2606.07348v1 (CC BY 4.0)}
\begin{document}
\maketitle

\noindent\emph{Source and licence.} Four intuitionistic modal connectives. Version arXiv:2606.07348v1 (CC BY 4.0), distributed by arXiv under the Creative Commons Attribution 4.0 International licence, http://creativecommons.org/licenses/by/4.0/, as recorded in the arXiv licence metadata for this exact version and shown on the abstract page https://arxiv.org/abs/2606.07348v1. Full licence notice: \texttt{licenses/arxiv-2606.07348-CC-BY-4.0.txt}.\par\smallskip

\noindent\textit{Editorial scope.} Counted unit: the article's proposition Proposition:let:be:a:balanced:clip:let:be:a:formula:if:then:for:all:balanced:clips:if:and:then:there:exists:a:balanced:clip:such:that:and:1:lozenge:Gamma:plus (source0.tex lines 4921--4959). The article's complete original proof of it is delivered with it (source0.tex lines 4962--5223, one \texttt{proof} environment of 262 lines). No mathematics is written, repaired or re-derived: the statement and the proof are the article's own lines, in the article's own notation, and no retained line is substituted or re-flowed.  No further source-local context is needed: every cross-reference of every retained line resolves inside the two delivered spans.  Nothing was omitted from any retained span, so the package carries no omission record.  No retained line contains a citation, so the wrapper carries no bibliography and no bibliography entry.  No retained line contains a figure environment, an image-inclusion command or drawing code, so no image is delivered and the package declares no asset dependency.  Editorial formatting only, in addition: an AMS \texttt{amsart} preamble that restates the article's own declarations and macros used by the retained lines, namely the environments \texttt{proposition}, \texttt{claim} on separate counters, together with the standard packages those lines need.  The retained environments print in this document as Proposition 1, Claim 1 in order of appearance, whereas the article prints the same items with the numbering of its own full document.  The complete native source of the article as distributed in the arXiv e-print for this exact version is bundled unchanged as \texttt{sources/202252/source0.tex}.
\begin{proposition}[Existence Property for $\lozenge$]\label{Proposition:let:be:a:balanced:clip:let:be:a:formula:if:then:for:all:balanced:clips:if:and:then:there:exists:a:balanced:clip:such:that:and:1:lozenge:Gamma:plus}
%
%
Let $((\Gamma^{+},\Gamma^{-}),A^{+},A^{-})$ be a balanced clip.
Let $D$ be a formula.
If ${\lozenge}D{\in}\Gamma^{+}$ then there exists a balanced clip $((\Delta^{+},\Delta^{-}),B^{+},
%
%
%
%
%
%
%
$\linebreak$
%
%
%
%
%
%
B^{-})$ such that $\Delta^{+}{\subseteq}\Gamma^{+}$ and $\Gamma^{-}{\subseteq}\Delta^{-}$ and there exists a balanced clip $((\Lambda^{+},\Lambda^{-}),C^{+},
%
%
%
%
%
%
%
$\linebreak$
%
%
%
%
%
%
C^{-})$ such that ${\square}\Delta^{+}{\subseteq}\Lambda^{+}$, ${\lozenge}\Delta^{-}{\subseteq}\Lambda^{-}$, $B^{-}{\leftharpoondown}B^{+}{\in}\Lambda^{+}$, $B^{+}{\rightharpoonup}B^{-}{\in}\Lambda^{-}$ and $D{\in}\Lambda^{+}$.
%
%
\end{proposition}

\begin{proof}\footnote{The reader is invited to see where is the only use of the inference rule $(\Rule8^{+})$ in this proof.}
%
%
Suppose ${\lozenge}D{\in}\Gamma^{+}$.
Let ${\mathcal S}$ be the set of all filters $\Lambda$ such that for all $E{\in}\Fo$, if $E{\in}\Lambda$ then ${\lozenge}E{\in}\Gamma^{+}$ and $D{\in}\Lambda$.
%
%
\begin{claim}
%
%
$\L^{+}{+}D$ is in ${\mathcal S}$.
%
%
\end{claim}
%
%
\begin{proof}
%
%
For the sake of the contradiction, suppose $\L^{+}{+}D$ is not in ${\mathcal S}$.
Hence, there exists $E{\in}\Fo$ such that $E{\in}\L^{+}{+}D$ and ${\lozenge}E{\not\in}\Gamma^{+}$.
Thus, $D{\rightharpoonup}E{\in}\L^{+}$.
Consequently, ${\lozenge}D{\rightharpoonup}{\lozenge}E{\in}\L^{+}$ using $(\Rule3^{+})$.
Since ${\lozenge}D{\in}\Gamma^{+}$, therefore ${\lozenge}E{\in}\Gamma^{+}$: a contradiction.
%
%
\medskip
%
%
\end{proof}
%
%
Moreover, for all nonempty chains $(\Lambda_{i})_{i{\in}I}$ in ${\mathcal S}$, $\bigcup\{\Lambda_{i}\ :\ i{\in}I\}$ is in ${\mathcal S}$.
Hence, by Kuratowski-Zorn Lemma, ${\mathcal S}$ possesses a maximal element $\Lambda$.
It is a routine task to show that $\Lambda$ is a prime filter such that for all $E{\in}\Fo$, if $E{\in}\Lambda$ then ${\lozenge}E{\in}\Gamma^{+}$ and $D{\in}\Lambda$.
Let ${\mathcal T}$ be the set of all filters $\Delta$ such that for all $G,H{\in}\Fo$, if $G{\vee}{\square}H{\in}\Delta$ then either $G{\in}\Gamma^{+}$, or $H{\in}\Lambda$ and for all $I{\in}\Fo$, if $I{\in}\Lambda$ then ${\lozenge}I{\in}\Delta$.
%
%
\begin{claim}
%
%
$\L^{+}{+}\{{\lozenge}E:\ E{\in}\Lambda\}$ is in ${\mathcal T}$.
%
%
\end{claim}
%
%
\begin{proof}
%
%
%This is the most difficult part of the proof of Prop.~\ref{Proposition:let:be:a:balanced:clip:let:be:a:formula:if:then:for:all:balanced:clips:if:and:then:there:exists:a:balanced:clip:such:that:and:1:lozenge:Gamma:plus}.
For the sake of the contradiction, suppose $\L^{+}{+}\{{\lozenge}E:\ E{\in}\Lambda\}$ is not in ${\mathcal T}$.
Hence, either there exists $G,H{\in}\Fo$ such that $G{\vee}{\square}H{\in}\L^{+}{+}\{{\lozenge}E:\ E{\in}\Lambda\}$, $G{\not\in}\Gamma^{+}$ and $H{\not\in}\Lambda$, or there exists $I{\in}\Fo$ such that $I{\in}\Lambda$ and ${\lozenge}I{\not\in}\L^{+}{+}\{{\lozenge}E:\ E{\in}\Lambda\}$.
%
%
\\
\\
%
%
In the former case, there exists $n{\in}\N$ and there exists $E_{1},\ldots,E_{n}{\in}\Lambda$ such that ${\lozenge}E_{1}{\wedge}\ldots
%
%
%
%
%
%
%
$\linebreak$
%
%
%
%
%
%
{\wedge}{\lozenge}E_{n}{\rightharpoonup}G{\vee}{\square}H{\in}\L^{+}$.
Since $H{\not\in}\Lambda$, therefore by the maximality of $\Lambda$ in ${\mathcal S}$, there exists $E^{\prime}{\in}\Fo$ such that $E^{\prime}{\in}\Lambda{+}H$ and ${\lozenge}E^{\prime}{\not\in}\Gamma^{+}$.
Thus, $H{\rightharpoonup}E^{\prime}{\in}\Lambda$.
Since $(E_{1}{\rightharpoonup}\ldots(E_{n}{\rightharpoonup}
%
%
%
%
%
%
%
$\linebreak$
%
%
%
%
%
%
((H{\rightharpoonup}E^{\prime}){\rightharpoonup}E_{1}{\wedge}\ldots{\wedge}E_{n}{\wedge}(H{\rightharpoonup}E^{\prime})))\ldots){\in}\L^{+}$ using $(\Formula18^{+})$ and $E_{1},\ldots,E_{n}{\in}\Lambda$,
%
%
%
%
%
%
%
\linebreak
%
%
%
%
%
%
therefore $E_{1}{\wedge}\ldots{\wedge}E_{n}{\wedge}(H{\rightharpoonup}E^{\prime}){\in}\Lambda$.
Consequently, ${\lozenge}(E_{1}{\wedge}\ldots{\wedge}E_{n}{\wedge}(H{\rightharpoonup}E^{\prime})){\in}
%
%
%
%
%
%
%
$\linebreak$
%
%
%
%
%
%
\Gamma^{+}$.
Since $({\lozenge}(E_{1}{\wedge}\ldots{\wedge}E_{n}{\wedge}(H{\rightharpoonup}E^{\prime})){\rightharpoonup}{\lozenge}E_{1}{\wedge}\ldots{\wedge}{\lozenge}E_{n}){\rightharpoonup}(({\lozenge}E_{1}{\wedge}\ldots{\wedge}{\lozenge}E_{n}{\rightharpoonup}G
%
%
%
%
%
%
%
$\linebreak$
%
%
%
%
%
%
{\vee}{\square}H){\rightharpoonup}({\lozenge}(E_{1}{\wedge}\ldots{\wedge}E_{n}{\wedge}(H{\rightharpoonup}E^{\prime})){\rightharpoonup}G{\vee}{\square}H)){\in}\L^{+}$ using $(\Formula16^{+})$, ${\lozenge}(E_{1}{\wedge}\ldots{\wedge}
%
%
%
%
%
%
%
$\linebreak$
%
%
%
%
%
%
E_{n}{\wedge}(H{\rightharpoonup}E^{\prime})){\rightharpoonup}{\lozenge}E_{1}{\wedge}\ldots{\wedge}{\lozenge}E_{n}{\in}\L^{+}$ using $(\Formula17^{+})$ and ${\lozenge}E_{1}{\wedge}\ldots{\wedge}{\lozenge}E_{n}{\rightharpoonup}G{\vee}{\square}H
%
%
%
%
%
%
%
$\linebreak$
%
%
%
%
%
%
{\in}\L^{+}$, therefore ${\lozenge}(E_{1}{\wedge}\ldots{\wedge}E_{n}{\wedge}(H{\rightharpoonup}E^{\prime})){\rightharpoonup}G{\vee}{\square}H{\in}\L^{+}$ using $(\Rule1^{+})$.
Hence,
%
%
%
%
%
%
%
\linebreak$
%
%
%
%
%
%
{\lozenge}(E_{1}{\wedge}\ldots{\wedge}E_{n}{\wedge}(H{\rightharpoonup}E^{\prime})){\rightharpoonup}G{\vee}{\lozenge}(E_{1}{\wedge}\ldots{\wedge}E_{n}{\wedge}(H{\rightharpoonup}E^{\prime}){\wedge}H){\in}\L^{+}$ using $(\Rule8^{+})$.
Since ${\lozenge}(E_{1}{\wedge}\ldots{\wedge}E_{n}{\wedge}(H{\rightharpoonup}E^{\prime})){\in}\Gamma^{+}$, therefore $G{\vee}{\lozenge}(E_{1}{\wedge}\ldots{\wedge}E_{n}{\wedge}(H{\rightharpoonup}E^{\prime}){\wedge}H)
%
%
%
%
%
%
%
$\linebreak$
%
%
%
%
%
%
{\in}\Gamma^{+}$.
Thus, either $G{\in}\Gamma^{+}$, or ${\lozenge}(E_{1}{\wedge}\ldots{\wedge}E_{n}{\wedge}(H{\rightharpoonup}E^{\prime}){\wedge}H){\in}\Gamma^{+}$.
Since $G{\not\in}\Gamma^{+}$, therefore ${\lozenge}(E_{1}{\wedge}\ldots{\wedge}E_{n}{\wedge}(H{\rightharpoonup}E^{\prime}){\wedge}H){\in}\Gamma^{+}$.
Since $E_{1}{\wedge}\ldots{\wedge}E_{n}{\wedge}(H{\rightharpoonup}E^{\prime}){\wedge}H{\rightharpoonup}E^{\prime}{\in}
%
%
%
%
%
%
%
$\linebreak$
%
%
%
%
%
%
\L^{+}$ using $(\Formula19^{+})$, therefore ${\lozenge}(E_{1}{\wedge}\ldots{\wedge}E_{n}{\wedge}(H{\rightharpoonup}E^{\prime}){\wedge}H){\rightharpoonup}{\lozenge}E^{\prime}{\in}\L^{+}$ using $(\Rule3^{+})$.
Since ${\lozenge}(E_{1}{\wedge}\ldots{\wedge}E_{n}{\wedge}(H{\rightharpoonup}E^{\prime}){\wedge}H){\in}\Gamma^{+}$, therefore ${\lozenge}E^{\prime}{\in}\Gamma^{+}$ a contradiction.
%
%
\\
\\
%
%
In the latter case, since ${\lozenge}I{\rightharpoonup}{\lozenge}I{\in}\L^{+}$ using $(\Formula15^{+})$, therefore ${\lozenge}I{\in}\L^{+}{+}\{{\lozenge}E:\ E{\in}\Lambda\}$: a contradiction.
%
%
\medskip
%
%
\end{proof}
%
%
Moreover, for all nonempty chains $(\Delta_{i})_{i{\in}I}$ in ${\mathcal T}$, $\bigcup\{\Delta_{i}\ :\ i{\in}I\}$ is in ${\mathcal T}$.
Hence, by Kuratowski-Zorn Lemma, ${\mathcal T}$ possesses a maximal element $\Delta$.
It is a routine task to show that $\Delta$ is a prime filter such that for all $G,H{\in}\Fo$, if $G{\vee}{\square}H{\in}\Delta$ then either $G{\in}\Gamma^{+}$, or $H{\in}\Lambda$ and for all $I{\in}\Fo$, if $I{\in}\Lambda$ then ${\lozenge}I{\in}\Delta$.
And in the end, choose $\Delta^{+}{=}\Delta$, $\Delta^{-}{=}\Fo{\setminus}\Delta$, $B^{+}{=}{\top}$, $B^{-}{=}{\bot}$, $\Lambda^{+}{=}\Lambda$, $\Lambda^{-}{=}\Fo{\setminus}\Lambda$, $C^{+}{=}{\top}$ and $C^{-}{=}{\bot}$.
Thus, $\Delta^{+}{\subseteq}
%
%
%
%
%
%
%
$\linebreak$
%
%
%
%
%
%
\Gamma^{+}$, $\Gamma^{-}{\subseteq}\Delta^{-}$, ${\square}\Delta^{+}{\subseteq}\Lambda^{+}$, ${\lozenge}\Delta^{-}{\subseteq}\Lambda^{-}$, $B^{-}{\leftharpoondown}B^{+}{\in}\Lambda^{+}$, $B^{+}{\rightharpoonup}B^{-}{\in}\Lambda^{-}$ and $D{\in}\Lambda^{+}$.
%
%
\medskip
%
%
\end{proof}

\end{document}
